\color{unidarkgreen}{\subsection{Сигнализация}}
\color{black}

\begin{enumerate}

\item
Устройство формирует выходное воздействие на сигнализацию \textit{«Сигнализация: Срабатывание»} и \textit{«Сигнализация: Срабатывание фикс.»} по схеме «ИЛИ» в следующих случаях: 
\begin{itemize}
    \item при отключении БСК от логики отключения основных, резервных и режимных защит;
    \item при отключении БСК от внешних защит;     
    \item при срабатывании сигнальной ступени ЗПН, ЗВП и ЗПВГ;
    \item Срабатывании ЗНР; 
    \item Срабатывании УРОВ.
\end{itemize}

\item
Устройство формирует выходное воздействие на сигнализацию \textit{«Сигнализация: Неисправность»} и \textit{«Сигнализация: Неисправность фикс.»} по схеме «ИЛИ» в следующих случаях:
\begin{itemize}
    \item при неисправности токовых цепей дифференциальных защит; 
    \item при низкой изоляции ТТ БСК;
    \item при неисправности цепей напряжения; 
    \item при неисправностях выключателя.
\end{itemize}

\item
Алгоритм работы логики логики предупредительной сигнализации приведен на рисунке \ref{pict_ls:LS}.

\begin{figure}[H]
\centering
\includegraphics[width=0.75\textwidth,height=0.75\textheight,keepaspectratio]{\fbpath/09.01 БСК/91 Сигнализация/img/LS.pdf}
\captionof{figure}{Функциональная схема логики блока сигнализации}\label{pict_ls:LS}
\end{figure}

\end{enumerate} 